\chapter{АНАЛИЗ ПРЕДМЕТНОЙ ОБЛАСТИ И ПОСТАНОВКА ЗАДАЧИ}

\section{Роль расчёта термодинамических свойств воды в инженерных задачах}

Вода и водяной пар являются базовой рабочей средой для теплоэнергетики,
химико-технологических процессов и смежных инженерных систем.
На этапах проектирования, наладки и эксплуатации требуется многократный пересчёт
состояния по ограниченному набору измеряемых параметров.
К типовым задачам относятся расчёты тепловых схем, выбор режимов котлоагрегатов и турбин,
оценка потерь в трубопроводах, анализ теплообмена и переходных процессов.

Результат расчёта используется не только как набор чисел.
Он служит основой для диагностических диаграмм и принятия эксплуатационных решений.
Диаграммы $T$--$s$, $p$--$h$, $p$--$v$ позволяют контролировать положение точки
относительно линии насыщения и быстрее выявлять недопустимые режимы.

С инженерной точки зрения критичны три свойства вычислительного инструмента.
Первое свойство --- скорость, достаточная для интерактивной работы и пакетной обработки.
Второе свойство --- устойчивость на фазовых границах и в околокритической области.
Третье свойство --- диагностируемость: при ошибочном вводе система должна выдавать
явную причину отказа, а не неинтерпретируемый результат.

\section{Стандарты и существующие подходы к расчёту свойств воды и пара}

Исторически расчёты выполнялись по паровым таблицам и диаграммам с интерполяцией~\cite{alexandrov_grigorev_1999}.
Современные программные реализации опираются на аналитические зависимости
и дополняют их численными методами для обратных маршрутов.

В инженерной практике применяются стандартизованные формулировки IAPWS.
Они обеспечивают сопоставимость результатов, но требуют строгого соблюдения
границ применимости и правил перехода между областями состояния.
Ключевая особенность предметной области --- региональная структура модели.
Разные участки рабочей области описываются разными уравнениями состояния
и разными независимыми переменными.

Наиболее сложным с вычислительной точки зрения является регион~3.
Для этой области естественными переменными являются плотность и температура,
поэтому при других входных парах требуются обратные зависимости и численные решатели.
Следовательно, практическая программная реализация должна включать не только формулы,
но и маршрутизацию расчёта, валидацию входа и обработку граничных случаев.

\subsection{Эволюция стандартов: от IFC-67 к IAPWS-IF97}

Для промышленного применения свойства воды и водяного пара
долгое время рассчитывались по формулировке IFC-67,
которая была положена в основу изданий ASME Steam Tables 1967 года
и многочисленных национальных таблиц и справочников~\cite{iapws_faq_asme}.
Именно эта традиция табличного представления расчётных данных
сохранилась и в российской инженерной практике,
что подтверждается поздними русскоязычными справочными изданиями
и учебной литературой~\cite{alexandrov_grigorev_1999,gidaspov_severina_2014}.

Однако развитие вычислительной техники
и рост требований к численной скорости расчёта
сделали табличный подход недостаточным.
В конце 1997 года IAPWS приняла промышленную формулировку IAPWS-IF97,
которая официально заменила IFC-67
и стала международным стандартом для расчётов в паросиловой отрасли~\cite{iapws_faq_asme}.
Новая формулировка сохраняет инженерную направленность,
но строится уже как набор региональных аналитических уравнений,
ориентированных на быструю и устойчивую машинную обработку.

По сравнению с IAPWS-95,
которая предназначена для общего и научного применения,
IF97 сознательно оптимизирована под промышленную эксплуатацию:
она быстрее вычисляется,
дольше сохраняется без смены нормативной базы
и лучше соответствует задачам расчёта тепловых схем,
оборудования и режимов работы энергетических установок~\cite{iapws_faq_asme,coolprop_if97_docs}.
Для дипломной работы это важно по двум причинам.
Во-первых, сама постановка задачи относится к инженерным, а не к фундаментальным расчётам.
Во-вторых, требуется единое ядро,
пригодное как для пользовательских интерфейсов,
так и для сервисной пакетной обработки.

\subsection{Сравнительный анализ существующих библиотек}

Современная программная практика уже предлагает готовые библиотеки,
использующие IF97 или совместимые с ней маршруты расчёта.
Наиболее показательными для сравнения являются библиотека CoolProp
и специализированная библиотека SteamTables.jl.

CoolProp представляет собой зрелую многокомпонентную библиотеку,
ориентированную на широкий круг термодинамических расчётов.
Для воды и пара в ней предусмотрена отдельная реализация \texttt{IF97::Water},
основанная на последних релизах IAPWS по термодинамическим и транспортным свойствам~\cite{coolprop_if97_docs}.
Преимуществами CoolProp являются
широкий набор обёрток для разных языков,
большое число поддерживаемых выходных свойств
и наличие большинства типовых входных пар.
При этом IF97-режим в CoolProp остаётся водо-паровым частным случаем:
он ограничен более узкой областью применимости,
не поддерживает смеси
и не реализует универсальные частные производные в общем интерфейсе~\cite{coolprop_if97_docs}.

SteamTables.jl, напротив, является компактной предметной библиотекой,
сосредоточенной только на воде и водяном паре.
Она реализует свойства IF97,
поддерживает входные постановки $(p, T)$, $(p, h)$ и $(p, s)$,
а также предоставляет функции насыщения
и возможность использовать физические единицы через \texttt{Unitful.jl}~\cite{steamtables_jl}.
Такой подход удобен для расчётных сценариев в среде Julia,
но библиотека ориентирована на более узкий круг маршрутов
и на один язык программирования.

Сопоставление этих решений показывает,
что существующие библиотеки хорошо закрывают задачу одиночного вызова формул,
но не решают автоматически архитектурную задачу проекта:
единое вычислительное ядро на Rust,
общее DTO-представление для разных клиентов,
маршрут $\rho T$ на всей рабочей области,
контролируемую маршрутизацию по регионам
и повторное использование одного и того же ядра
в настольном интерфейсе, оболочке на базе веб-технологий и gRPC-сервисе.
Следовательно, в рамках дипломной работы требуется
не только выбрать стандарт расчёта,
но и реализовать программный комплекс,
в котором физическая модель, архитектура и эксплуатационные требования
согласованы между собой.

\section{Обоснование выбора IAPWS-IF97 как основы модели}

В работе принят стандарт IAPWS-IF97 (редакция 2012 года) как нормативная
и физическая основа вычислительной модели~\cite{iapws_if97_2012}.
Выбор обусловлен промышленной направленностью стандарта,
формализованной топологией области расчёта
и наличием обратных зависимостей для практических маршрутов.
По сравнению с табличными подходами и библиотеками общего назначения
IF97 даёт нормативно определённый,
достаточно быстрый
и инженерно интерпретируемый базис для построения собственного вычислительного ядра.

IF97 задаёт уравнения для регионов 1--5, линию насыщения (регион~4)
и правила перехода между областями, включая границу B23.
Это позволяет построить алгоритм, в котором регион определяется автоматически,
а вычисления выполняются по профильной модели выбранной области.

Стандарт не содержит единой универсальной аналитической формулы для всех входных пар
во всей рабочей области.
Для региона~3 используются специализированные обратные зависимости,
в том числе релиз для $v(p,T)$~\cite{iapws_vpt_region3}.
Маршрут $(\rho, T)$ реализуется как численное расширение.
Таким образом, IF97 остаётся неизменной физической базой,
а численные методы используются как инженерное дополнение,
не нарушающее нормативную интерпретацию результата.

\section{Требования к программному комплексу}

Требования формулируются исходя из инженерных сценариев и ограничений IF97.
Ключевой принцип: единый источник вычислений и одинаковый результат
во всех каналах доставки.

Функциональные требования:
\begin{itemize}
\item расчёт состояния воды и водяного пара в области применимости IAPWS-IF97;
\item поддержка маршрутов $pT$, $ph$, $ps$, $px$, $\rho T$ по входным параметрам $p$, $T$, $h$, $s$, $\rho$, $x$;
\item автоматическое определение региона IF97 с учётом линии насыщения и границы B23;
\item возврат унифицированного набора выходных свойств ($v$, $\rho$, $h$, $s$, $u$, $c_p$, $w$) и номера региона;
\item валидация диапазонов и корректная диагностика ошибок ввода и расчёта;
\item единый контракт обмена через \texttt{if97\_app\_api} для пользовательских приложений и сервисного контура.
\end{itemize}

Нефункциональные требования:
\begin{itemize}
\item вычислительная эффективность для интерактивного режима и пакетных расчётов;
\item воспроизводимость сборки и детерминированность результата;
\item устойчивость на фазовых границах и в околокритической области с явной сигнализацией о несходимости;
\item отсутствие дублирования предметной математики в пользовательских и сервисных слоях;
\item возможность масштабирования пакетной обработки и контроля нагрузки;
\item расширяемость архитектуры за счёт разделения ядра и тонких адаптерных слоёв;
\item сохранение физической интерпретации IF97 при численных расширениях (в частности, для $\rho T$).
\end{itemize}

Архитектура реализуется по схеме «одно вычислительное ядро и несколько каналов доставки».
Ядро \texttt{if97\_core} выполняет расчёт и валидацию.
Слой \texttt{if97\_app\_api} фиксирует DTO и контракты обмена.
Поверх ядра работают FLTK-клиент,
клиент на базе Yew, WebAssembly и Tauri
и gRPC-сервис,
что исключает расхождение формул между компонентами.

В качестве эксплуатационного ориентира для сервисного канала принимается
пропускная способность порядка $10^5$ точек/с
в пакетных gRPC-сценариях.
Актуальный контрольный прогон
подтверждает данный порядок величины:
среднее значение составило 129~632,83~точек/с.
Этот ориентир используется как проверяемый критерий при развитии системы.

\section{Постановка вычислительной задачи и поддерживаемые режимы ввода-вывода}

Вычислительная задача формулируется как отображение
из допустимой входной пары параметров в полное состояние рабочей среды.
Множество входных параметров: $p$, $T$, $h$, $s$, $\rho$, $x$.
Множество выходных параметров: $v$, $\rho$, $h$, $s$, $u$, $c_p$, $w$, номер региона IF97.

Для каждой допустимой входной пары требуется:
\begin{itemize}
\item определить регион IF97 и корректный маршрут расчёта;
\item вычислить унифицированный набор выходных свойств;
\item вернуть диагностируемую ошибку при недопустимом вводе или выходе за границы применимости.
\end{itemize}

Отдельное требование относится к региону~4.
Система должна корректно обрабатывать линию насыщения и двухфазные состояния,
поскольку эти режимы напрямую связаны с практическими задачами кипения,
конденсации и оценки влажности пара.

Поддерживаемые режимы ввода-вывода (маршруты расчёта) включают:
\begin{itemize}
\item прямой маршрут $pT$: расчёт состояния по давлению и температуре;
\item обратные маршруты $ph$ и $ps$: расчёт по давлению и энтальпии или энтропии;
\item маршрут $px$: расчёт по давлению и степени сухости в двухфазной области;
\item маршрут $\rho T$: численное расширение на основе метода секущих
с подбором давления и последующим вычислением по $pT$.
\end{itemize}

Во всех режимах обязательны валидация входа и контроль области применимости.
При выходе за границы IF97 или при некорректной комбинации параметров
возвращается диагностируемая ошибка, а не неопределённые значения.

Критерием корректности постановки является
совпадение физической модели и формата результата
для всех каналов доставки
(\texttt{if97\_core}, FLTK,
связки Yew, WebAssembly и Tauri,
gRPC).
Критерием эксплуатационной применимости является
сохранение целевого порядка производительности
в пакетных сценариях и в контейнерном запуске.

Сформулированная постановка задачи определяет логику следующих глав:
математическую модель и численные методы (глава 2),
реализацию \texttt{if97\_core} (глава 3),
клиентские приложения (глава 4)
и сервисно-инфраструктурный контур (глава 5).
